% SGA 3, Expose VI_A, Loop 1 English reconstruction.
% Scope: section 4 heading and introductory paragraph, followed by
% paragraph 4.1 in full; paragraph 4.2 is excluded.
% Authority: Polo--Gille Exp6A-13oct24.pdf, component-local pages 26--27
% (combined-reader pages 329--330; printed pages 318--319).
% OCR and the Reinhold English Markdown are comparison witnesses only.

\section*{4. 商的构造 \(F\backslash G\)
  (一般情形)}
\phantomsection\label{sec:sga3-VIA-4}
\addcontentsline{toc}{section}{4. 商的构造
  \texorpdfstring{\(F\backslash G\)}{F backslash G}
  (一般情形)}

现在采用\hyperref[thm:sga3-VIA-3.2]{定理~3.2} 的假设,
但不要求 \(F\) 和 \(G\) 在 \(A\) 上有限型.

\medskip
\phantomsection\label{par:sga3-VIA-4.1}%
\noindent\textbf{4.1.---}
设 \(G^\alpha\) 为 \(G\) 的连通分支. 先证明 \(G^\alpha\)
对群胚 \(G_*\) 定义的等价关系的饱和
\[
  \mathscr S(G^\alpha)
\]%
\footnote{\label{note:sga3-VIA-44}编者注: 我们用 \(\mathscr S(G^\alpha)\) 而不是
\(\overline{G^\alpha}\) 表示 \(G^\alpha\) 的饱和.}
为 \(G\) 的既开又闭的子集; 换言之, 它是 \(G\) 的若干连通分支之并.

此饱和是 \(F\times_A G^\alpha\) 在 \(\lambda\) 下的像,
因而在 \(G\) 中开 (参见 \S~\hyperref[par:sga3-VIA-3.1]{3.1}).
设 \(k\) 为 \(A\) 的剩余域, \(\bar k\) 为 \(k\) 的代数闭包.
只需证明
\[
  (F\times_A G^\alpha)\otimes_A k
\]
在 \(\lambda\otimes_A k\) 下的像于 \(G\otimes_A k\) 中闭;
或由 SGA~1, VIII.4.4, 等价地证明
\[
  (F\times_A G^\alpha)\otimes_A\bar k
\]
在 \(\lambda\otimes_A\bar k\) 下的像于 \(G\otimes_A\bar k\) 中闭.
由于 \(G^\alpha\otimes_A\bar k\) 为 \(G\otimes_A\bar k\) 的有限多个连通分支之并,
可以归结为 \(A\) 为代数闭域的情形, 此后便作此假设.
此时 \(\mathscr S(G^\alpha)\) 为 \(G^\alpha\) 在左平移
\(\ell_{u(x)}\) 下的像之并, 其中 \(x\) 遍历 \(F\) 的闭点.
由于这些像为 \(G\) 的连通分支, 所述断言成立.
